\scanpage{8}
\uline{Corollaire} Soit $S$ sur $\mathbb F_q$. La classification des
torseurs $P$ sous $G_S$, munis de $P\simeq P^f$, i.e
de $\varphi:P\mathbin{\text{---}}P$ satisfaisant $\varphi(xg)=\varphi(x)f(g)$, équivaut
à celle des \deleted{rev. \inserted{étales} \unsure{galoisiens} de $S$, de groupe}
$G(\mathbb F_q)$.

En effet, $N$ est le gr. discret \inserted{sur $k$} (constant)
défini par $G(\mathbb F_q)$, et on applique le n° 1.

\textcircled{3}. Soit $S$ préschéma en car. $p$, et $E$ un
Module (qu. coh.) sur $S$. Une $p$-structure sur $E$
est un hom
\[
\varphi:E^{(p)}\longrightarrow E
\]
[où $E^{(p)}=E\otimes_{\underline O_S}(\underline O_S,f)$]. La $p$-structure est
dite \uline{non dégénérée} si $\varphi$ est un isom.
On veut classifier les $p$-structures non
dégénérées sur $S$, avec $E$ loc. libre de rang $n$.
Ce sont les structures de $G$-torseurs avec $f$-
rigidification, où $G=\mathrm{GL}(n)_S$, et $f:G\mathbin{\text{---}}G$
est le frobenius. Comme $G$ provient de
$G_0$ sur $\mathbb F_p$, on peut appliquer 2°, on trouve que
\mspage{7}
